%  LaTeX support: latex@mdpi.com 
%  For support, please attach all files needed for compiling as well as the log file, and specify your operating system, LaTeX version, and LaTeX editor.

%=================================================================
\documentclass[metlit,article,accept,pdftex,moreauthors]{Definitions/mdpi} 
\usepackage{caption}
\usepackage{subcaption}
\usepackage[indonesian]{babel}
% \linenumbers
%--------------------

%---------
% pdftex
%---------
% The option pdftex is for use with pdfLaTeX. Remove "pdftex" for (1) compiling with LaTeX & dvi2pdf (if eps figures are used) or for (2) compiling with XeLaTeX.

%=================================================================
% MDPI internal commands - do not modify
\firstpage{1} 
\makeatletter 
\setcounter{page}{\@firstpage} 
\makeatother
\pubvolume{1}
\issuenum{1}
\articlenumber{0}
\pubyear{2024}
\copyrightyear{2024}
%\externaleditor{Academic Editor: Firstname Lastname}
\datereceived{ } 
\daterevised{ } % Comment out if no revised date
\dateaccepted{ } 
\datepublished{ } 
%\datecorrected{} % For corrected papers: "Corrected: XXX" date in the original paper.
%\dateretracted{} % For corrected papers: "Retracted: XXX" date in the original paper.
\hreflink{https://doi.org/} % If needed use \linebreak
%\doinum{}
%\pdfoutput=1 % Uncommented for upload to arXiv.org
%\CorrStatement{yes}  % For updates


%=================================================================
% Add packages and commands here. The following packages are loaded in our class file: fontenc, inputenc, calc, indentfirst, fancyhdr, graphicx, epstopdf, lastpage, ifthen, float, amsmath, amssymb, lineno, setspace, enumitem, mathpazo, booktabs, titlesec, etoolbox, tabto, xcolor, colortbl, soul, multirow, microtype, tikz, totcount, changepage, attrib, upgreek, array, tabularx, pbox, ragged2e, tocloft, marginnote, marginfix, enotez, amsthm, natbib, hyperref, cleveref, scrextend, url, geometry, newfloat, caption, draftwatermark, seqsplit
% cleveref: load \crefname definitions after \begin{document}

%=================================================================
% Please use the following mathematics environments: Theorem, Lemma, Corollary, Proposition, Characterization, Property, Problem, Example, ExamplesandDefinitions, Hypothesis, Remark, Definition, Notation, Assumption
%% For proofs, please use the proof environment (the amsthm package is loaded by the MDPI class).

%=================================================================
% Full title of the paper (Capitalized)
\Title{Judul Penelitian Mahasiswa Program Sarjana Teknologi Informasi}

% MDPI internal command: Title for citation in the left column
\TitleCitation{Judul Penelitian Mahasiswa}

% Author Orchid ID: enter ID or remove command
\newcommand{\orcidauthorA}{0000-0000-0000-000X} % Add \orcidA{} behind the author's name
%\newcommand{\orcidauthorB}{0000-0000-0000-000X} % Add \orcidB{} behind the author's name

% Authors, for the paper (add full first names)
\Author{Joko Subianto$^{1}$, Prabowo Widodo$^{2}$}

%\longauthorlist{yes}

% MDPI internal command: Authors, for metadata in PDF
\AuthorNames{Joko Subianto dan Prabowo Widodo}

% MDPI internal command: Authors, for citation in the left column
\AuthorCitation{Subianto, J.; Widodo, P.}
% If this is a Chicago style journal: Lastname, Firstname, Firstname Lastname, and Firstname Lastname.

% Affiliations / Addresses (Add [1] after \address if there is only one affiliation.)
\address{%
$^{1}$ \quad Teknologi Informasi, Telkom University, Bandung, Indonesia\\
$^{2}$ \quad Informatika, Telkom University, Bandung, Indonesia\\
}

% Contact information of the corresponding author
\corres{Correspondence: jokop@telkomuniversity.ac.id; Tel.: 0812 (eks; 233) +62-813-322-7887 (F.L.)}

% Current address and/or shared authorship
\firstnote{Alamat saat ini: Bandung, Dayeuhkolot.}  % Current address should not be the same as any items in the Affiliation section.
\secondnote{Semua penulis telah berkontribusi.}
% The commands \thirdnote{} till \eighthnote{} are available for further notes

%\simplesumm{} % Simple summary

%\conference{} % An extended version of a conference paper

% Abstract (Do not insert blank lines, i.e. \\) 
\abstract{Ini adalah tugas menulis artikel jurnal. Abstrak harus berisi antara lain ringkasan masalah, ringkasan metode, dan ringkasan hasil. Abstrak sebaiknya ditulis dalam satu paragraf. Jika penelitian kuantitatif sebutkan beberapa angka yang menunjukkan hasil pengukuran.
}

% Keywords
\keyword{Kata Kunci; Abstrak; Metode Penelitian } 

% The fields PACS, MSC, and JEL may be left empty or commented out if not applicable
%\PACS{J0101}
%\MSC{}
%\JEL{}

%%%%%%%%%%%%%%%%%%%%%%%%%%%%%%%%%%%%%%%%%%
% Only for the journal Diversity
%\LSID{\url{http://}}

%%%%%%%%%%%%%%%%%%%%%%%%%%%%%%%%%%%%%%%%%%
% Only for the journal Applied Sciences
%\featuredapplication{Authors are encouraged to provide a concise description of the specific application or a potential application of the work. This section is not mandatory.}
%%%%%%%%%%%%%%%%%%%%%%%%%%%%%%%%%%%%%%%%%%

%%%%%%%%%%%%%%%%%%%%%%%%%%%%%%%%%%%%%%%%%%
% Only for the journal Data
%\dataset{DOI number or link to the deposited data set if the data set is published separately. If the data set shall be published as a supplement to this paper, this field will be filled by the journal editors. In this case, please submit the data set as a supplement.}
%\datasetlicense{License under which the data set is made available (CC0, CC-BY, CC-BY-SA, CC-BY-NC, etc.)}

%%%%%%%%%%%%%%%%%%%%%%%%%%%%%%%%%%%%%%%%%%
% Only for the journal Toxins
%\keycontribution{The breakthroughs or highlights of the manuscript. Authors can write one or two sentences to describe the most important part of the paper.}

%%%%%%%%%%%%%%%%%%%%%%%%%%%%%%%%%%%%%%%%%%
% Only for the journal Encyclopedia
%\encyclopediadef{For entry manuscripts only: please provide a brief overview of the entry title instead of an abstract.}

%%%%%%%%%%%%%%%%%%%%%%%%%%%%%%%%%%%%%%%%%%
% Only for the journal Advances in Respiratory Medicine
%\addhighlights{yes}
%\renewcommand{\addhighlights}{%

%\noindent This is an obligatory section in “Advances in Respiratory Medicine”, whose goal is to increase the discoverability and readability of the article via search engines and other scholars. Highlights should not be a copy of the abstract, but a simple text allowing the reader to quickly and simplified find out what the article is about and what can be cited from it. Each of these parts should be devoted up to 2~bullet points.\vspace{3pt}\\
%\textbf{What are the main findings?}
% \begin{itemize}[labelsep=2.5mm,topsep=-3pt]
% \item First bullet.
% \item Second bullet.
% \end{itemize}\vspace{3pt}
%\textbf{What is the implication of the main finding?}
% \begin{itemize}[labelsep=2.5mm,topsep=-3pt]
% \item First bullet.
% \item Second bullet.
% \end{itemize}
%}

%%%%%%%%%%%%%%%%%%%%%%%%%%%%%%%%%%%%%%%%%%
\begin{document}

%%%%%%%%%%%%%%%%%%%%%%%%%%%%%%%%%%%%%%%%%%
\section{Pendahuluan}
Pendahuluan memberikan masalah yang harus disajikan dan diperkuat dengan data atau statistik seperti nominal USD 27.79 Milyar pada 2025. Singkatan yang baru pertama kali muncul harus ditulis definisinya seperti ini Compound Annual Growth Rate (CAGR). Selanjutnya apabila ada singkatan yang sama, singkatan bisa disebut secara langsung seperti CAGR. Sitasi menggunakan format IEEE \cite{berita2024}. Referensi lebih dari satu pada satu kalimat bisa menggunakan contoh berikut \cite{Nuha2021,IndustrialWSN2023}. Semua referensi harus dipastikan sudah disitasi.

Jika Penulis telah banyak membaca penelitian sebelumnya, penulis bisa menuliskan celah penelitian yang belum dilakukan sebelumnya dan akan dijawab pada penelitian ini. Penulis harus menuliskan pertanyaan penelitian pada bagian pendahuluan ini. Selain itu, penulis harus menuliskan hipotesis pada bagian ini. Sebaiknya, penulis mengakhiri bagian ini dengan struktur penulisan.




\section{Penelitian Terkait}
Bagian ini berisi penelitian terkait. Kumpulkan semua karya ilmiah terkait. Setiap artikel harus disitasi dan ditampilkan di referensi. Nama penulis bisa disebut seperti ini Silva Fre dkk \cite{Silva2015}. Baca semua artikel terkait lalu kelompokkan dan buat menjadi ringkasan-ringkasan dalam paragraf.

\section{Metodologi}
Metodologi sebaiknya digambarkan menggunakan diagram dan juga rumus jika diperlukan. Metodologi ini harus menjelaskan cara mengambil kesimpulan. Selain itu, jika terdapat metode atau algoritma baru, bisa ditampilkan juga.

\subsection{Menampilkan Rumus}
Persamaan sering juga disebut rumus. Formula juga menunjukkan persamaan. Persamaan berikut bisa digunakan untuk menghitung daya,
\begin{equation}
\label{eq:friis}
P_R = \frac{A_r A_t}{d^2 \lambda^2} P_t    
\end{equation}


Pastikan setiap variabel pada rumus sudah dijelaskan. Seperti, $P_R$ adalah daya dan $\lambda$ adalah panjang gelombang. Selain itu rumus ini juga bisa digunakan:
\begin{equation}\label{eq:derivative_f_t}
f'(t) = \frac{d}{dt} \left[ \cos\left(\frac{\pi t}{2T}\right) - \frac{T - t}{T} \right].
\end{equation}


\subsection{Menampilkan Gambar}

Setiap gambar yang muncul harus disebutkan di dalam teks seperti Gambar \ref{fig:comparison a} . Selain gambar, bisa juga ditampilkan potongan code yang disebut listing. Metode juga bisa diilustraikan menggunakan pseudocode.
\begin{figure}[H]
    \centering
    \includegraphics[width=0.6\linewidth]{image_rev/diagram_flowchart.png}
    \caption{Diagram Alir Sistem \cite{xia2018path}}
    \label{fig:comparison a}
\end{figure}

Selain gambar dan tabel, artikel bisa ditambahkan daftar. Sebaiknya daftar disertai penomoran namun dengan item pun tidak mengapa. Pembuatan daftar bisa menggunakan item :
\begin{itemize}
    \item \text{ Hal pertama (\( D \))}: 30
    \item \text{ Hal kedua (\( N \))}: 30
    \item \text{ Hal ketiga (\( T \))}: 100
    \item \text{ Hal keempat }: 30
\end{itemize}

\subsection{Menampilkan Tabel}


Selain gambar, setiap tabel harus disebutkan dalam teks seperti Tabel \ref{tab:gwo_comparison}. Kalimat ini adalah kalimat kedua. Sedangkan ini adalah kalimat ke tiga.

\begin{table}[H]
\centering
\caption{Perbandingan Data.}
\begin{tabular}{@{}lcccc@{}}
\toprule
Cara & \multicolumn{1}{c}{ekspektasi} & \multicolumn{1}{c}{varian} & \multicolumn{1}{p{2.8cm}}{\centering rat. \\  iterasi} & \multicolumn{1}{p{2.8cm}}{\centering Waktu \\ ekse (s)} \\ \midrule
 garis      & \( 0.12 \) & \( 0.003 \) & 23 & 9.2 \\
kuadrat   & \( 0.13 \) & \( 0.003 \) & 23 & 9 \\
eks & \( 0.003 \) & \( 0.0034 \) & 45 & 0.02 \\
kos      & \( 0.002 \) & \( 0.023 \) & 67 & 0.250 \\ 
\bottomrule
\end{tabular}
\label{tab:gwo_comparison}
\end{table}




\vspace{5pt}
\section{Hasil Percobaan}
Percobaan dapat diringkas dalam bentuk tabel seperti pada Tabel \ref{tab:params}. Gambar digunakan untuk menunjukkan sampel data dan pola atau kecenderungan hasil seperti pada Gambar \ref{fig:pso_vs_simplistic}. Kalimat ini adalah kalimat terakhir.





\begin{table}[H]
\centering
\caption{Data Hasil Percobaan}
\label{tab:params}
\begin{tabular}{@{}lcc@{}}
\toprule
aspek  & Nilai [Satuan] \\ \midrule
angka angka penting & 20 \\
data 1 & 4 \\ data 2 & 34 x 31 \\
A & 0.88 \\
B & 40 \\
(kecil;  besar) & [0, 20] [kg] \\
berat badan & 50 \\
% Add more rows as needed
\bottomrule
\end{tabular}
\end{table}

\begin{figure}[H] 
    \centering
    \includegraphics[width=0.8\textwidth]{image_rev/gambar_hasil.png}
    \caption{Perbandingan hasil \cite{xia2018path}}
    \label{fig:pso_vs_simplistic}
\end{figure}




 Setiap gambar dan tabel harus dibahas dalam teks. Cara membahas data dalam tabel dengan menunjukkan nilai yang menonjol seperti nilai tertinggi maupun terendah. Sedangkan grafik menunjukkan kecenderungan atau pola, seperti pola naik atau turun. Ini adalah penutup paragraf.



\section{Simpulan}

Simpulan sama seperti abstrak yang mana meringkas hasil percobaan dan kesimpulan yang didapatkan. Ini adalah lanjutan kalimat sebelumnya. Ini adalah kalimat terakhir.

%%%%%%%%%%%%%%%%%%%%%%%%%%%%%%%%%%%%%%%%%%
\vspace{6pt} 

%%%%%%%%%%%%%%%%%%%%%%%%%%%%%%%%%%%%%%%%%%
%% optional
%\supplementary{The following supporting information can be downloaded at:  \linksupplementary{s1}, Figure S1: title; Table S1: title; Video S1: title.}

% Only for journal Methods and Protocols:
% If you wish to submit a video article, please do so with any other supplementary material.
% \supplementary{The following supporting information can be downloaded at: \linksupplementary{s1}, Figure S1: title; Table S1: title; Video S1: title. A supporting video article is available at doi: link.}

% Only for journal Hardware:
% If you wish to submit a video article, please do so with any other supplementary material.
% \supplementary{The following supporting information can be downloaded at: \linksupplementary{s1}, Figure S1: title; Table S1: title; Video S1: title.\vspace{6pt}\\
%\begin{tabularx}{\textwidth}{lll}
%\toprule
%\textbf{Name} & \textbf{Type} & \textbf{Description} \\
%\midrule
%S1 & Python script (.py) & Script of python source code used in XX \\
%S2 & Text (.txt) & Script of modelling code used to make Figure X \\
%S3 & Text (.txt) & Raw data from experiment X \\
%S4 & Video (.mp4) & Video demonstrating the hardware in use \\
%... & ... & ... \\
%\bottomrule
%\end{tabularx}
%}

%%%%%%%%%%%%%%%%%%%%%%%%%%%%%%%%%%%%%%%%%%
%\authorcontributions{For research articles with several authors, a short paragraph specifying their individual contributions must be provided. The following statements should be used ``Conceptualization, X.X. and Y.Y.; methodology, X.X.; software, X.X.; validation, X.X., Y.Y. and Z.Z.; formal analysis, X.X.; investigation, X.X.; resources, X.X.; data curation, X.X.; writing---original draft preparation, X.X.; writing---review and editing, X.X.; visualization, X.X.; supervision, X.X.; project administration, X.X.; funding acquisition, Y.Y. All authors have read and agreed to the published version of the manuscript.'', please turn to the  \href{http://img.mdpi.org/data/contributor-role-instruction.pdf}{CRediT taxonomy} for the term explanation. Authorship must be limited to those who have contributed substantially to the work~reported.}
%
%\funding{Please add: ``This research received no external funding'' or ``This research was funded by NAME OF FUNDER grant number XXX.'' and  and ``The APC was funded by XXX''. Check carefully that the details given are accurate and use the standard spelling of funding agency names at \url{https://search.crossref.org/funding}, any errors may affect your future funding.}
%
%\institutionalreview{In this section, you should add the Institutional Review Board Statement and approval number, if relevant to your study. You might choose to exclude this statement if the study did not require ethical approval. Please note that the Editorial Office might ask you for further information. Please add “The study was conducted in accordance with the Declaration of Helsinki, and approved by the Institutional Review Board (or Ethics Committee) of NAME OF INSTITUTE (protocol code XXX and date of approval).” for studies involving humans. OR “The animal study protocol was approved by the Institutional Review Board (or Ethics Committee) of NAME OF INSTITUTE (protocol code XXX and date of approval).” for studies involving animals. OR “Ethical review and approval were waived for this study due to REASON (please provide a detailed justification).” OR “Not applicable” for studies not involving humans or animals.}
%
%\informedconsent{Any research article describing a study involving humans should contain this statement. Please add ``Informed consent was obtained from all subjects involved in the study.'' OR ``Patient consent was waived due to REASON (please provide a detailed justification).'' OR ``Not applicable'' for studies not involving humans. You might also choose to exclude this statement if the study did not involve humans.
%
%Written informed consent for publication must be obtained from participating patients who can be identified (including by the patients themselves). Please state ``Written informed consent has been obtained from the patient(s) to publish this paper'' if applicable.}
%
%\dataavailability{We encourage all authors of articles published in MDPI journals to share their research data. In this section, please provide details regarding where data supporting reported results can be found, including links to publicly archived datasets analyzed or generated during the study. Where no new data were created, or where data is unavailable due to privacy or ethical restrictions, a statement is still required. Suggested Data Availability Statements are available in section ``MDPI Research Data Policies'' at \url{https://www.mdpi.com/ethics}.} 

% Only for journal Nursing Reports
%\publicinvolvement{Please describe how the public (patients, consumers, carers) were involved in the research. Consider reporting against the GRIPP2 (Guidance for Reporting Involvement of Patients and the Public) checklist. If the public were not involved in any aspect of the research add: ``No public involvement in any aspect of this research''.}

% Only for journal Nursing Reports
%\guidelinesstandards{Please add a statement indicating which reporting guideline was used when drafting the report. For example, ``This manuscript was drafted against the XXX (the full name of reporting guidelines and citation) for XXX (type of research) research''. A complete list of reporting guidelines can be accessed via the equator network: \url{https://www.equator-network.org/}.}

% Only for journal Nursing Reports
%\useofartificialintelligence{Please describe in detail any and all uses of artificial intelligence (AI) or AI-assisted tools used in the preparation of the manuscript. This may include, but is not limited to, language translation, language editing and grammar, or generating text. Alternatively, please state that “AI or AI-assisted tools were not used in drafting any aspect of this manuscript”.}

\acknowledgments{Penulis mengucapkan terima kasih pada dosen kelas. Data yang digunakan pada penelitian ini dapat diakses pada https://github.com/hilalnuha/EGWO.}


\conflictsofinterest{Penulis menyatakan tidak ada konflik kepentingan.} 

\useofartificialintelligence{Penulisan dibantu oleh AI menggunakan DeepSeek}


%%%%%%%%%%%%%%%%%%%%%%%%%%%%%%%%%%%%%%%%%%
%% Optional

%% Only for journal Encyclopedia
%\entrylink{The Link to this entry published on the encyclopedia platform.}
%
%\abbreviations{Abbreviations}{
%The following abbreviations are used in this manuscript:\\
%
%\noindent 
%\begin{tabular}{@{}ll}
%MDPI & Multidisciplinary Digital Publishing Institute\\
%DOAJ & Directory of open access journals\\
%TLA & Three letter acronym\\
%LD & Linear dichroism
%\end{tabular}
%}
%
%%%%%%%%%%%%%%%%%%%%%%%%%%%%%%%%%%%%%%%%%%%
%%% Optional
%\appendixtitles{no} % Leave argument "no" if all appendix headings stay EMPTY (then no dot is printed after "Appendix A"). If the appendix sections contain a heading then change the argument to "yes".
%\appendixstart
%\appendix
%\section[\appendixname~\thesection]{}
%\subsection[\appendixname~\thesubsection]{}
%The appendix is an optional section that can contain details and data supplemental to the main text---for example, explanations of experimental details that would disrupt the flow of the main text but nonetheless remain crucial to understanding and reproducing the research shown; figures of replicates for experiments of which representative data are shown in the main text can be added here if brief, or as Supplementary Data. Mathematical proofs of results not central to the paper can be added as an appendix.
%
%\begin{table}[H] 
%\caption{This is a table caption.\label{tab5}}
%\newcolumntype{C}{>{\centering\arraybackslash}X}
%\begin{tabularx}{\textwidth}{CCC}
%\toprule
%\textbf{Title 1}	& \textbf{Title 2}	& \textbf{Title 3}\\
%\midrule
%Entry 1		& Data			& Data\\
%Entry 2		& Data			& Data\\
%\bottomrule
%\end{tabularx}
%\end{table}
%
%\section[\appendixname~\thesection]{}
%All appendix sections must be cited in the main text. In the appendices, Figures, Tables, etc. should be labeled, starting with ``A''---e.g., Figure A1, Figure A2, etc.

%%%%%%%%%%%%%%%%%%%%%%%%%%%%%%%%%%%%%%%%%%
\begin{adjustwidth}{-\extralength}{0cm}
%\printendnotes[custom] % Un-comment to print a list of endnotes

\reftitle{Referensi}

% Please provide either the correct journal abbreviation (e.g. according to the “List of Title Word Abbreviations” http://www.issn.org/services/online-services/access-to-the-ltwa/) or the full name of the journal.
% Citations and References in Supplementary files are permitted provided that they also appear in the reference list here. 

%=====================================
% References, variant A: external bibliography
%=====================================
\bibliography{sumber}

%=====================================
% References, variant B: internal bibliography
%=====================================


% If authors have biography, please use the format below
%\section*{Short Biography of Authors}
%\bio
%{\raisebox{-0.35cm}{\includegraphics[width=3.5cm,height=5.3cm,clip,keepaspectratio]{Definitions/author1.pdf}}}
%{\textbf{Firstname Lastname} Biography of first author}
%
%\bio
%{\raisebox{-0.35cm}{\includegraphics[width=3.5cm,height=5.3cm,clip,keepaspectratio]{Definitions/author2.jpg}}}
%{\textbf{Firstname Lastname} Biography of second author}

% For the MDPI journals use author-date citation, please follow the formatting guidelines on http://www.mdpi.com/authors/references
% To cite two works by the same author: \citeauthor{ref-journal-1a} (\citeyear{ref-journal-1a}, \citeyear{ref-journal-1b}). This produces: Whittaker (1967, 1975)
% To cite two works by the same author with specific pages: \citeauthor{ref-journal-3a} (\citeyear{ref-journal-3a}, p. 328; \citeyear{ref-journal-3b}, p.475). This produces: Wong (1999, p. 328; 2000, p. 475)

%%%%%%%%%%%%%%%%%%%%%%%%%%%%%%%%%%%%%%%%%%
%% for journal Sci
%\reviewreports{\\
%Reviewer 1 comments and authors’ response\\
%Reviewer 2 comments and authors’ response\\
%Reviewer 3 comments and authors’ response
%}
%%%%%%%%%%%%%%%%%%%%%%%%%%%%%%%%%%%%%%%%%%
\PublishersNote{}
\end{adjustwidth}
\end{document}

